\bmtchapitre{Progression indicative}{Organiser la dépendance des notions et adapter les séances au calendrier de l’établissement.}{Édition du professeur}
Les quatre parties du livre sont des regroupements éditoriaux. Elles n’imposent pas de commencer par l’analyse. La référence PE T11 conservée au dépôt attribue $65$ h à sa composante Analyse (avec équations, arithmétique et boucles), $20$ h à Algèbre (logique et suites), $45$ h à Géométrie et $20$ h aux données et probabilités : $150$ h au total. Les heures d’évaluation et de remédiation se répartissent dans ces enveloppes.
\begin{center}\begin{tabular}{L{2cm}L{9cm}}\toprule
Séquence&Chapitres et dépendances\\\midrule
1&Logique (\ref{t11s-c06}), boucles (\ref{t11s-c07}), équations et systèmes (\ref{t11s-c08}–\ref{t11s-c09}).\\
2&Arithmétique et bases (\ref{t11s-c10}–\ref{t11s-c11}); vecteurs et produit scalaire (\ref{t11s-c13}), puis trigonométrie (\ref{t11s-c16}).\\
3&Fonctions, limites, branches, dérivation et étude (\ref{t11s-c01}–\ref{t11s-c05}).\\
4&Suites (\ref{t11s-c12}); barycentres et transformations (\ref{t11s-c14}–\ref{t11s-c15}).\\
5&Dénombrement (\ref{t11s-c19}), espace (\ref{t11s-c17}), probabilités (\ref{t11s-c20}) et statistique (\ref{t11s-c18}).\\\bottomrule
\end{tabular}\end{center}
La RAPE archivée concerne 2025–2026; les dates de l’année 2026–2027 ne sont pas reprises sans document actuel vérifié. Introduire les fonctions trigonométriques après les rappels du cercle. Faire les DS de partie quand tous leurs prérequis sont installés, même si les chapitres ont été enseignés dans un autre ordre.
\section{Remédiation}
Après une erreur, identifier sa nature : calcul, domaine, choix de méthode, justification ou lecture. Reprendre un exemple court puis demander un exercice de structure différente. Pour une classe nombreuse, comparer deux solutions au tableau et faire justifier la différence; éviter une correction réduite au résultat final.
